\subsubsection{詳細設計}

\term{詳細設計}{Detailed Design} は、主要構成要素の内部を実装できる水準まで具体化する設計である。これは、プログラマがコードを書くための直接的な手がかりになる。

詳細設計では、主に次を扱う。
\begin{itemize}
  \item 主要構成要素を内部モジュールやプログラム単位へ分ける。
  \item 各モジュールの責務を割り当てる。
  \item モジュール間のやり取りを決める。
  \item 名前の有効範囲、可視性、公開・非公開の境界を決める。
  \item 状態、モード、状態遷移を定義する。
  \item データと制御の依存関係を明らかにする。
  \item データ構造、アルゴリズム、ユーザインタフェースの詳細を決める。
\end{itemize}

詳細設計は、実装と完全に分離して行われるとは限らない。特に反復型やアジャイル型の開発では、設計と実装が近い間隔で進む。しかし、その場合でも、設計判断が存在しなくなるわけではない。設計判断がコードの中だけに閉じ込められるのか、明示的に記録されるのかが変わるだけである。
